\chapter{Orthogonality and Orthogonal Subspaces}

A central concept throughout linear algebra is \textbf{orthogonality}---the
mathematical generalization of perpendicularity and right angles. When two
directions are orthogonal, their components decouple completely: knowledge of
one coordinate gives zero information about the other. Because of this decoupling,
complex multi-dimensional problems separate into independent one-dimensional
subproblems.

Here we extend perpendicularity from pairs of individual vectors to whole
\emph{subspaces} -- lines, planes, and higher-dimensional subspaces that meet
at right angles. The four fundamental subspaces of a matrix partition the domain
and codomain into two orthogonal pairs, providing the structural foundation
for projections, least squares, and the singular value decomposition. We also
introduce \emph{orthonormal} sets -- coordinate systems composed of mutually
perpendicular unit vectors -- and the matrices built from them, the
\emph{orthogonal matrices}, which represent rotations and reflections. These
transformations preserve lengths and angles exactly, making them stable tools
in numerical linear algebra.

As always: picture first, then algebra, then proof.

\section{The dot product remembers the angle}

Everything here rides on one number. Given two vectors
$\vec{u},\vec{v}\in\R^n$ written as columns, their \emph{dot product} (or inner product) is
\[
  \inner{\vec{u}}{\vec{v}} \;=\; \vec{u}\T\vec{v}
  \;=\; u_1 v_1 + u_2 v_2 + \cdots + u_n v_n.
\]
It is a sum of products, and the answer is a single scalar. The length of a vector is recovered
by dotting it with itself:
\[
  \norm{\vec{v}}^2 \;=\; \inner{\vec{v}}{\vec{v}} \;=\; v_1^2 + \cdots + v_n^2,
\]
which is just the Pythagorean theorem in $n$ dimensions. And the dot product remembers the
\emph{angle} $\theta$ between the two vectors through the formula
\[
  \inner{\vec{u}}{\vec{v}} \;=\; \norm{\vec{u}}\,\norm{\vec{v}}\,\cos\theta.
\]

\begin{intuitionbox}
The dot product $\inner{\vec{u}}{\vec{v}} = \norm{\vec{u}}\,\norm{\vec{v}}\,\cos\theta$
measures the directional alignment between two vectors. When two nonzero vectors point in
perpendicular directions, $\cos\theta = 0$, so their dot product vanishes:
$\inner{\vec{u}}{\vec{v}} = 0$. In linear algebra, vectors with zero inner product are called
\textbf{orthogonal}.
\end{intuitionbox}

We know the angle formula from Chapter~4, but it pays to see \emph{why} a zero dot product means a
right angle from first principles, because the geometric reason is pure Pythagoras and needs no
trigonometry at all.

\begin{figure}[H]
\centering
\begin{tikzpicture}[scale=1.2,>=Stealth]
  \coordinate (O) at (0,0);
  \coordinate (U) at (2.6,0);
  \coordinate (V) at (0,1.8);
  \draw[->,spacegray] (-0.4,0)--(3.2,0);
  \draw[->,spacegray] (0,-0.4)--(0,2.4);
  \draw[->,vecblue,very thick] (O)--(U) node[below right]{$\vec{u}$};
  \draw[->,vecred,very thick]  (O)--(V) node[above left]{$\vec{v}$};
  \draw[->,structurepurple,very thick] (V)--(U);
  \node[structurepurple,above right] at (1.3,0.9) {$\vec{u}-\vec{v}$};
  % right angle mark
  \draw[spacegray] (0,0.25)--(0.25,0.25)--(0.25,0);
\end{tikzpicture}
\caption{When $\vec{u}\perp\vec{v}$, the triangle with legs $\vec{u}$ and $\vec{v}$ and
hypotenuse $\vec{u}-\vec{v}$ obeys Pythagoras: $\norm{\vec{u}-\vec{v}}^2=\norm{\vec{u}}^2+\norm{\vec{v}}^2$.
That equation is \emph{equivalent} to $\inner{\vec{u}}{\vec{v}}=0$.}
\label{ch17:pythagoras}
\end{figure}

\begin{formalbox}
\textbf{Orthogonality via Pythagoras.} For any $\vec{u},\vec{v}\in\R^n$,
\[
  \norm{\vec{u}-\vec{v}}^2
  = \inner{\vec{u}-\vec{v}}{\vec{u}-\vec{v}}
  = \norm{\vec{u}}^2 - 2\inner{\vec{u}}{\vec{v}} + \norm{\vec{v}}^2.
\]
The Pythagorean relation $\norm{\vec{u}-\vec{v}}^2 = \norm{\vec{u}}^2 + \norm{\vec{v}}^2$ holds
\emph{if and only if} the cross term vanishes, that is, if and only if
$\inner{\vec{u}}{\vec{v}}=0$. So ``right angle'' and ``zero dot product'' are two names for the
same thing.
\end{formalbox}

\begin{recoverybox}
\textbf{Pause and think: The diagonals of a rectangle.} If $\vec{u}\perp\vec{v}$, what is the
squared length of their sum, $\norm{\vec{u}+\vec{v}}^2$? Expand it:
\[
  \norm{\vec{u}+\vec{v}}^2 = \norm{\vec{u}}^2 + 2\inner{\vec{u}}{\vec{v}} + \norm{\vec{v}}^2
  = \norm{\vec{u}}^2 + \norm{\vec{v}}^2.
\]
Because the cross term $2\inner{\vec{u}}{\vec{v}}$ is zero, $\norm{\vec{u}+\vec{v}}^2$ equals
$\norm{\vec{u}-\vec{v}}^2$. Geometrically, $\vec{u}+\vec{v}$ and $\vec{u}-\vec{v}$ are the two
diagonals of the parallelogram spanned by $\vec{u}$ and $\vec{v}$. When the angle is $90^\circ$,
that parallelogram is a rectangle, and its two diagonals have the exact same length.
\end{recoverybox}

That little derivation is the model for the whole chapter: expand a squared length, watch the
cross terms, and demand that they vanish. Let us make the definition official.

\begin{definition}\label{ch17:def-orthogonal-vectors}
Two vectors $\vec{u},\vec{v}\in\R^n$ are \textbf{orthogonal} if $\inner{\vec{u}}{\vec{v}}
= \vec{u}\T\vec{v} = 0$. We write $\vec{u}\perp\vec{v}$.
\end{definition}

\begin{remark}
By this definition, the zero vector is orthogonal to every vector $\vec{v}\in\R^n$,
since $\inner{\vec{0}}{\vec{v}} = 0$. This convention ensures that orthogonal
complements are linear subspaces containing $\vec{0}$.
\end{remark}

\begin{example}\label{ch17:first-numeric}
Let $\vec{u}=\colvec{3\\1\\-2}$ and $\vec{v}=\colvec{2\\2\\4}$. Then
\[
  \inner{\vec{u}}{\vec{v}} = (3)(2)+(1)(2)+(-2)(4) = 6+2-8 = 0,
\]
so $\vec{u}\perp\vec{v}$. They sit at a right angle in $\R^3$, even though neither one is
along an axis. Check Pythagoras as a sanity test: $\norm{\vec{u}}^2 = 9+1+4 = 14$,
$\norm{\vec{v}}^2 = 4+4+16 = 24$, and $\vec{u}-\vec{v}=\colvec{1\\-1\\-6}$ has
$\norm{\vec{u}-\vec{v}}^2 = 1+1+36 = 38 = 14+24$. The cross term is zero.
\end{example}

\begin{trapbox}
Orthogonality must not be confused with collinearity. Nonzero parallel vectors
have maximal inner product in absolute value; orthogonal vectors have an inner
product of zero. In addition, while every set of nonzero mutually orthogonal
vectors is linearly independent (as proved below), linearly independent vectors
need not be orthogonal.
\end{trapbox}

\section{From vectors to subspaces}

Now we lift the idea. So far ``orthogonal'' relates two arrows. The reach grows once we
ask a whole \emph{line} to be perpendicular to a whole \emph{plane}, or one subspace to another.

\begin{definition}\label{ch17:def-orthogonal-subspaces}
Two subspaces $V$ and $W$ of $\R^n$ are \textbf{orthogonal}, written $V\perp W$, if
\emph{every} vector in $V$ is orthogonal to \emph{every} vector in $W$:
\[
  \inner{\vec{v}}{\vec{w}} = 0 \qquad \text{for all } \vec{v}\in V,\ \vec{w}\in W.
\]
\end{definition}

This is a strong demand -- it must hold for all pairs at once -- but it is much easier to check
than it looks. A small lemma reduces the work to testing \emph{basis} vectors.

\begin{lemma}\label{ch17:basis-suffices}
Let $V=\spn(\vec{v}_1,\dots,\vec{v}_p)$ and $W=\spn(\vec{w}_1,\dots,\vec{w}_q)$. If every basis
vector of $V$ is orthogonal to every basis vector of $W$ -- that is, $\inner{\vec{v}_i}{\vec{w}_j}=0$
for all $i,j$ -- then $V\perp W$.
\end{lemma}

\begin{proof}
Take any $\vec{v}\in V$ and $\vec{w}\in W$. Write them in the bases: $\vec{v}=\sum_i a_i\vec{v}_i$
and $\vec{w}=\sum_j b_j\vec{w}_j$. The dot product is bilinear (it distributes over sums and pulls
out scalars in each slot), so
\[
  \inner{\vec{v}}{\vec{w}}
  = \Big\langle \textstyle\sum_i a_i\vec{v}_i,\ \sum_j b_j\vec{w}_j\Big\rangle
  = \sum_i\sum_j a_i b_j \,\inner{\vec{v}_i}{\vec{w}_j}
  = \sum_i\sum_j a_i b_j \cdot 0 = 0. \qedhere
\]
\end{proof}

\begin{example}\label{ch17:line-perp-plane}
In $\R^3$, let $V$ be the $z$-axis, spanned by $\vec{e}_3=\colvec{0\\0\\1}$, and let $W$ be the
$xy$-plane, spanned by $\vec{e}_1=\colvec{1\\0\\0}$ and $\vec{e}_2=\colvec{0\\1\\0}$. Each basis
dot product is zero: $\inner{\vec{e}_3}{\vec{e}_1}=0$ and $\inner{\vec{e}_3}{\vec{e}_2}=0$. By the
lemma, the whole $z$-axis is orthogonal to the whole $xy$-plane. Good -- that matches the picture
in your head.
\end{example}

\begin{trapbox}
Two \emph{planes} in $\R^3$ can never be orthogonal subspaces in the sense of
Definition~\ref{ch17:def-orthogonal-subspaces}, even though we casually call the floor and a wall
``perpendicular.'' Why? Two planes through the origin in $\R^3$ always share a whole line (their
intersection). Any nonzero vector $\vec{x}$ on that shared line lives in \emph{both} planes, so we
would need $\inner{\vec{x}}{\vec{x}}=0$, forcing $\vec{x}=\vec{0}$. The everyday ``perpendicular
walls'' meet along an edge that belongs to both -- so as \emph{subspaces} they are not orthogonal.
Orthogonal subspaces can share \emph{only} the zero vector.
\end{trapbox}

\begin{figure}[H]
\centering
\begin{tikzpicture}[scale=1.1,>=Stealth]
  % Plane V (horizontal floor)
  \fill[sourceblue!15] (-2.2,-0.9) -- (1.8,-0.9) -- (2.6,0.7) -- (-1.4,0.7) -- cycle;
  \draw[sourceblue!70,thick] (-2.2,-0.9) -- (1.8,-0.9) -- (2.6,0.7) -- (-1.4,0.7) -- cycle;
  \node[sourceblue!80!black,font=\small\bfseries] at (-1.4,0.4) {Plane $V$};
  
  % Plane W (vertical wall along the intersection line)
  \fill[quotientorange!15] (-0.5,-1.0) -- (0.5,1.0) -- (0.5,2.6) -- (-0.5,0.6) -- cycle;
  \draw[quotientorange!80,thick] (-0.5,-1.0) -- (0.5,1.0) -- (0.5,2.6) -- (-0.5,0.6) -- cycle;
  \node[quotientorange!90!black,font=\small\bfseries] at (0.6,2.4) {Plane $W$};

  % Shared intersection line L = V \cap W through the origin
  \draw[spacegray,dashed] (-0.8,-1.6) -- (0.8,1.6);
  \draw[vecred,line width=1.3pt] (-0.5,-1.0) -- (0.5,1.0);
  \fill[black] (0,0) circle (1.3pt) node[below right,font=\scriptsize] {$\vec{0}$};

  % A nonzero vector x lying on the shared line
  \draw[->,vecred,very thick] (0,0) -- (0.4,0.8) node[above left,font=\small] {$\vec{x}\in V\cap W$};
\end{tikzpicture}
\caption{Why two planes in $\R^3$ can never be orthogonal subspaces. Even if they meet at
$90^\circ$ like the floor and a wall, they intersect along an entire line $V\cap W$. Any nonzero
vector $\vec{x}$ on that line belongs to \emph{both} planes. If $V\perp W$, we would have
$\inner{\vec{x}}{\vec{x}}=0$, forcing $\vec{x}=\vec{0}$. True orthogonal subspaces can share
nothing except the origin.}
\label{ch17:planes-trap}
\end{figure}

That last observation is so important it deserves to be a proposition.

\begin{proposition}\label{ch17:trivial-intersection}
If $V\perp W$ then $V\cap W = \{\vec{0}\}$.
\end{proposition}

\begin{proof}
Suppose $\vec{x}\in V\cap W$. Then $\vec{x}$ is in $V$ and also in $W$, so by orthogonality
$\inner{\vec{x}}{\vec{x}}=0$. But $\inner{\vec{x}}{\vec{x}}=\norm{\vec{x}}^2$, and the only vector
with zero length is $\vec{0}$. Hence $\vec{x}=\vec{0}$.
\end{proof}

\begin{intuitionbox}
Orthogonal subspaces intersect only at the origin: $V \cap W = \{\vec{0}\}$.
Geometrically, every direction in $V$ is perpendicular to every direction in $W$,
as seen with the $z$-axis and the $xy$-plane in $\R^3$.
\end{intuitionbox}

\section{The orthogonal complement}

A line through the origin in $\R^3$ is orthogonal to many things, but one space perpendicular to it
is the largest possible: the entire plane at right angles. That ``everything perpendicular to $V$''
is useful enough to earn its own name and symbol.

\begin{definition}\label{ch17:def-complement}
The \textbf{orthogonal complement} of a subspace $V\subseteq\R^n$, written $V^{\perp}$ (read
``$V$ perp''), is the set of all vectors orthogonal to every vector in $V$:
\[
  V^{\perp} \;=\; \set{\,\vec{x}\in\R^n : \inner{\vec{x}}{\vec{v}}=0 \text{ for all }\vec{v}\in V\,}.
\]
\end{definition}

A few facts are immediate and worth pinning down, since $V^\perp$ is going to do a lot of work.

\begin{proposition}\label{ch17:perp-is-subspace}
For any subspace $V\subseteq\R^n$, the set $V^{\perp}$ is itself a subspace of $\R^n$.
\end{proposition}

\begin{proof}
We check the two closure rules. First, $\vec{0}\in V^\perp$ since $\inner{\vec{0}}{\vec{v}}=0$ for
all $\vec{v}$. Now suppose $\vec{x},\vec{y}\in V^\perp$ and $c$ is a scalar. For any $\vec{v}\in V$,
\[
  \inner{\vec{x}+\vec{y}}{\vec{v}} = \inner{\vec{x}}{\vec{v}} + \inner{\vec{y}}{\vec{v}} = 0+0 = 0,
  \qquad
  \inner{c\vec{x}}{\vec{v}} = c\,\inner{\vec{x}}{\vec{v}} = c\cdot 0 = 0.
\]
So $\vec{x}+\vec{y}$ and $c\vec{x}$ are again in $V^\perp$. Closed under addition and scaling: it
is a subspace.
\end{proof}

\begin{focusbox}
The two facts to carry with you about $V^\perp$:
\begin{itemize}
  \item $V^\perp$ is a subspace, and $V\perp V^\perp$ (they are orthogonal to each other).
  \item Together they fill the whole space cleanly: $\dimn V + \dimn V^\perp = n$, and
        $V\cap V^\perp=\{\vec0\}$. We say $\R^n = V \oplus V^\perp$.
\end{itemize}
Here the direct sum symbol $\oplus$ means that every vector in $\R^n$ decomposes into a piece in
$V$ plus a piece in $V^\perp$ in exactly one unique way, with zero overlap between the two
subspaces except the origin. ``Complement'' is the right word: $V^\perp$ supplies exactly the
dimensions $V$ is missing, with no overlap and at right angles.
\end{focusbox}

The dimension count $\dimn V + \dimn V^\perp = n$ comes shortly, once we connect $V^\perp$
to a nullspace -- the cleanest route to it. First, a hands-on computation to make the
idea concrete.

\begin{labbox}
Let $V=\spn\!\left(\colvec{1\\1\\1}\right)$, a line in $\R^3$. To find $V^\perp$ we want all
$\vec{x}=\colvec{x\\y\\z}$ with $\inner{\vec{x}}{(1,1,1)}=0$, i.e.\ $x+y+z=0$. That single equation
describes a \emph{plane} through the origin. A basis for it: set the free variables and solve.
Taking $(y,z)$ free gives $\colvec{-1\\1\\0}$ and $\colvec{-1\\0\\1}$. So $V^\perp$ is the
2-dimensional plane $x+y+z=0$. Check: we have $\dimn V=1$, $\dimn V^\perp=2$, and
$1+2=3=n$. The dimensions add up.
\end{labbox}

\begin{figure}[H]
\centering
\begin{tikzpicture}[scale=1.1,>=Stealth,line cap=round]
  % Plane V^\perp (drawn as a horizontal parallelogram in perspective)
  \fill[targetgreen!15] (-2.4,-0.8) -- (1.6,-0.8) -- (2.4,0.8) -- (-1.6,0.8) -- cycle;
  \draw[targetgreen!60,thick] (-2.4,-0.8) -- (1.6,-0.8) -- (2.4,0.8) -- (-1.6,0.8) -- cycle;
  \node[targetgreen!70!black,font=\small\bfseries] at (-1.5,0.5) {$V^{\perp}$ (plane)};
  % Line V: dashed below the plane, solid above
  \draw[vecblue,thick,dashed] (0,-1.5) -- (0,0);
  \fill[black] (0,0) circle (1.3pt) node[below left,font=\scriptsize] {$\vec{0}$};
  \draw[->,vecblue,very thick] (0,0) -- (0,1.9) node[above,font=\small\bfseries]{$V$ (line)};
  % Sample vector w in the plane
  \draw[->,targetgreen!70!black,very thick] (0,0) -- (1.7,0.3) node[right,font=\small] {$\vec{w}\in V^{\perp}$};
  % Right-angle mark between vertical line and plane vector
  \draw[spacegray] (0,0.3) -- (0.24,0.34) -- (0.24,0.04);
\end{tikzpicture}
\caption{A line $V$ and its orthogonal complement $V^{\perp}$, a plane in $\R^3$. Every vector $\vec{w}$
in the plane meets the line at $90^\circ$. Together their dimensions add to the full space
($1+2=3$), and they intersect only at the origin $\vec{0}$.}
\label{ch17:complement-pic}
\end{figure}

Notice the trick in the lab. Finding $V^\perp$ for a line turned into solving a single
linear equation, and that equation was exactly ``row vector dotted with $\vec x$ equals zero.''
Stack several such equations into the rows of a matrix and $V^\perp$ becomes a nullspace.

\section{The key identity: \texorpdfstring{$\nullsp(\mat A)=\row(\mat A)^\perp$}{null A = row A perp}}

One observation ties orthogonal complements to matrices and sets up the whole
``big picture.'' Look at what it means for $\vec{x}$ to be in the nullspace of a matrix
$\mat{A}$. The equation $\mat{A}\vec{x}=\vec{0}$, written row by row, says
\[
  (\text{row 1 of }\mat{A})\cdot\vec{x}=0,\quad
  (\text{row 2 of }\mat{A})\cdot\vec{x}=0,\quad\dots,\quad
  (\text{row }m\text{ of }\mat{A})\cdot\vec{x}=0.
\]
Every row of $\mat{A}$ dots to zero with $\vec{x}$. So $\vec{x}$ is orthogonal to every row -- and
by Lemma~\ref{ch17:basis-suffices}, orthogonal to every \emph{combination} of rows, i.e.\ to the
entire row space. That is the definition of being in $\row(\mat{A})^\perp$.

\begin{theorem}\label{ch17:null-is-row-perp}
For any $m\times n$ matrix $\mat{A}$,
\[
  \nullsp(\mat{A}) \;=\; \row(\mat{A})^{\perp},
\]
as subspaces of the input space $\R^n$. The nullspace and the row space are orthogonal
complements.
\end{theorem}

\begin{proof}
We show each set contains the other.

($\subseteq$) If $\vec{x}\in\nullsp(\mat{A})$ then $\mat{A}\vec{x}=\vec{0}$, so each row dotted
with $\vec{x}$ is zero. Hence $\vec{x}$ is orthogonal to every row, and by
Lemma~\ref{ch17:basis-suffices} to all of $\row(\mat{A})$. Thus $\vec{x}\in\row(\mat{A})^\perp$.

($\supseteq$) Conversely, if $\vec{x}\in\row(\mat{A})^\perp$ then $\vec{x}$ is orthogonal to every
row in particular, so each row dotted with $\vec{x}$ is zero, which says $\mat{A}\vec{x}=\vec{0}$.
Thus $\vec{x}\in\nullsp(\mat{A})$.

The two inclusions give equality.
\end{proof}

This theorem pays for the dimension law promised earlier. The Rank--Nullity theorem (from the
rank chapter) says $\dimn\row(\mat A)+\dimn\nullsp(\mat A)=n$, because the row space has dimension
$r$ (the rank) and the nullspace has dimension $n-r$. Read through Theorem~\ref{ch17:null-is-row-perp},
that becomes a statement purely about a subspace and its complement.

\begin{theorem}[Dimension of the complement]\label{ch17:dim-complement}
For any subspace $V\subseteq\R^n$,
\[
  \dimn V + \dimn V^{\perp} = n.
\]
\end{theorem}

\begin{proof}
Pick a basis $\vec{v}_1,\dots,\vec{v}_r$ of $V$ and stack these as the \emph{rows} of an
$r\times n$ matrix $\mat{A}$. Then $V=\row(\mat{A})$ and $\dimn V=r=\rank(\mat A)$. By
Theorem~\ref{ch17:null-is-row-perp}, $V^\perp=\row(\mat A)^\perp=\nullsp(\mat A)$, whose dimension
is $n-r$ by Rank--Nullity. So $\dimn V+\dimn V^\perp = r+(n-r)=n$.
\end{proof}

\begin{corollary}\label{ch17:perp-perp}
$(V^\perp)^\perp = V$. Taking the complement twice returns you to where you started.
\end{corollary}

\begin{proof}
Every vector of $V$ is orthogonal to every vector of $V^\perp$, so $V\subseteq (V^\perp)^\perp$.
By Theorem~\ref{ch17:dim-complement} applied twice, $\dimn (V^\perp)^\perp = n-\dimn V^\perp
= n-(n-\dimn V)=\dimn V$. A subspace contained in another of the same dimension must equal it.
\end{proof}

\begin{strangbox}
Strang reaches this point and draws his famous ``big picture'': two perpendicular pairs of
subspaces, the input space $\R^n$ split by the row space and the nullspace, the output space
$\R^m$ split by the column space and the left nullspace. \emph{Before} Strang, hold onto the
single sentence ``the nullspace is everything the rows can't see,'' which is exactly
$\nullsp(\mat A)=\row(\mat A)^\perp$. \emph{After} Strang, recognize this same orthogonality as the
reason projection onto the column space (next chapter) and least squares both work: the error
vector lands in the left nullspace, perpendicular to the column space.
\end{strangbox}

\section{The two orthogonal pairs among the four subspaces}

We now have the first orthogonal pair for free. Apply the identity to $\mat{A}\T$ instead of
$\mat{A}$, and the second pair follows, because the row space of $\mat{A}\T$ is the column space
of $\mat{A}$, and the nullspace of $\mat{A}\T$ is the left nullspace of $\mat A$.

\begin{theorem}[The two orthogonal pairs]\label{ch17:two-pairs}
Let $\mat{A}$ be $m\times n$ with rank $r$. Then:
\begin{itemize}
  \item In the input space $\R^n$:\quad
        $\nullsp(\mat A) = \row(\mat A)^{\perp}$,\quad and these have dimensions $n-r$ and $r$.
  \item In the output space $\R^m$:\quad
        $\nullsp(\mat A\T) = \col(\mat A)^{\perp}$,\quad and these have dimensions $m-r$ and $r$.
\end{itemize}
Consequently $\R^n = \row(\mat A)\oplus\nullsp(\mat A)$ and
$\R^m = \col(\mat A)\oplus\nullsp(\mat A\T)$, each a clean perpendicular split.
\end{theorem}

\begin{proof}
The first line is Theorem~\ref{ch17:null-is-row-perp}. For the second, replace $\mat A$ by
$\mat A\T$ in that theorem: $\nullsp(\mat A\T)=\row(\mat A\T)^\perp$. But the rows of $\mat A\T$
are the columns of $\mat A$, so $\row(\mat A\T)=\col(\mat A)$. Substituting gives
$\nullsp(\mat A\T)=\col(\mat A)^\perp$. The dimension counts come from
Theorem~\ref{ch17:dim-complement} and $\rank(\mat A\T)=\rank(\mat A)=r$.
\end{proof}

\begin{figure}[H]
\centering
\begin{tikzpicture}[scale=1.0,>=Stealth,
  spc/.style={draw=spacegray,thick,rounded corners=2pt},
  sub/.style={font=\small}]
  % Input space R^n on the left
  \begin{scope}[shift={(0,0)}]
    \node[spc,minimum width=3.6cm,minimum height=4.8cm] (Rn) at (0,0) {};
    \node[above=1pt of Rn,font=\bfseries] {$\R^{n}$ (inputs)};
    % row space (upper) and nullspace (lower) meeting at origin
    \fill[sourceblue!18] (-1.5,0.2) rectangle (1.5,2.2);
    \fill[quotientorange!20] (-1.5,-2.2) rectangle (1.5,-0.2);
    \node[sub,sourceblue!60!black] at (0,1.2) {$\row(\mat A)$};
    \node[sub,sourceblue!60!black] at (0,0.75) {$\dim r$};
    \node[sub,quotientorange!65!black] at (0,-1.1) {$\nullsp(\mat A)$};
    \node[sub,quotientorange!65!black] at (0,-1.55) {$\dim n-r$};
    \fill[black] (0,0) circle (1.5pt) node[right,font=\scriptsize] {$\vec{0}$};
    \node[spacegray,font=\small\bfseries] at (-0.35,0) {$\perp$};
  \end{scope}
  % arrow A
  \draw[->,very thick,structurepurple] (2.1,0) -- node[above]{$\mat A$} (4.5,0);
  % Output space R^m on the right
  \begin{scope}[shift={(6.6,0)}]
    \node[spc,minimum width=3.6cm,minimum height=4.8cm] (Rm) at (0,0) {};
    \node[above=1pt of Rm,font=\bfseries] {$\R^{m}$ (outputs)};
    \fill[targetgreen!20] (-1.5,0.2) rectangle (1.5,2.2);
    \fill[trapred!18] (-1.5,-2.2) rectangle (1.5,-0.2);
    \node[sub,targetgreen!55!black] at (0,1.2) {$\col(\mat A)$};
    \node[sub,targetgreen!55!black] at (0,0.75) {$\dim r$};
    \node[sub,trapred!75!black] at (0,-1.05) {$\nullsp(\mat A\T)$};
    \node[sub,trapred!75!black] at (0,-1.5) {$\dim m-r$};
    \fill[black] (0,0) circle (1.5pt) node[right,font=\scriptsize] {$\vec{0}$};
    \node[spacegray,font=\small\bfseries] at (-0.35,0) {$\perp$};
  \end{scope}
\end{tikzpicture}
\caption{The big picture. Inside $\R^n$ the row space and the nullspace are orthogonal complements;
inside $\R^m$ the column space and the left nullspace are orthogonal complements. The matrix
$\mat A$ maps the row space one-to-one onto the column space, and crushes the nullspace to the
origin. The little $\perp$ symbols remind you each pair meets at $90^\circ$ at the origin.}
\label{ch17:big-picture}
\end{figure}

\begin{aestheticbox}
A matrix is, on its face, just a grid of numbers.
Yet hidden inside every such grid are four subspaces, locked into two perpendicular pairs, with
dimensions $r$, $n-r$, $r$, $m-r$ that always add correctly. The rank $r$ appears \emph{twice} --
once on each side -- which is the deep reason ``row rank equals column rank.'' Orthogonality is
what makes the picture rigid: the complement isn't just some other subspace of the right
dimension, it is the \emph{unique} one at right angles. That rigidity is what lets us later
\emph{split} any vector uniquely into a part the matrix sees and a part it ignores -- the
projection behind least squares, regression, and the SVD.
\end{aestheticbox}

\begin{example}\label{ch17:four-spaces-worked}
Let
\[
  \mat A=\begin{bmatrix} 1 & 2 & 2\\ 2 & 4 & 5 \end{bmatrix}.
\]
Row reduce: $R_2\to R_2-2R_1$ gives $\begin{bmatrix}1&2&2\\0&0&1\end{bmatrix}$, then clear above the
second pivot: $R_1\to R_1-2R_2$ gives $\begin{bmatrix}1&2&0\\0&0&1\end{bmatrix}$. So $r=2$.

\textbf{Row space} ($\subseteq\R^3$, dimension $2$): spanned by the pivot rows
$\colvec{1\\2\\0}$ and $\colvec{0\\0\\1}$.

\textbf{Nullspace} ($\subseteq\R^3$, dimension $1$): the free variable is $x_2$. From
$x_1+2x_2=0$ and $x_3=0$ we get $\vec n=\colvec{-2\\1\\0}$. Check the orthogonality promised by
Theorem~\ref{ch17:two-pairs}: $\inner{\vec n}{(1,2,0)}=-2+2+0=0$ and $\inner{\vec n}{(0,0,1)}=0$.
The nullspace really is perpendicular to the row space.

\textbf{Column space} ($\subseteq\R^2$, dimension $2$): the pivots sit in columns 1 and 3, so a
basis is $\colvec{1\\2}$ and $\colvec{2\\5}$. These two independent vectors already fill $\R^2$, so
$\col(\mat A)=\R^2$.

\textbf{Left nullspace} ($\subseteq\R^2$, dimension $m-r=0$): just $\{\vec 0\}$, consistent with
the column space being all of $\R^2$ (its complement has nothing left to be). Everything checks.
\end{example}

\begin{walkbox}
Take a walk and answer this without paper. A matrix $\mat A$ is $4\times 6$ with rank $3$. What are
the dimensions of all four subspaces, and which two live in $\R^6$, which two in $\R^4$? Then:
which pair sums to $6$, which pair sums to $4$? If you can recite ``$3$ and $3$ in $\R^6$, $3$ and
$1$ in $\R^4$'' and say which is the row space and which the nullspace, the big picture has become
yours.
\end{walkbox}

\section{Orthonormal sets and coordinate expansions}

Orthogonality immediately implies linear independence.

\begin{proposition}\label{ch17:orthog-implies-indep}
A set of nonzero, mutually orthogonal vectors $\vec{q}_1,\dots,\vec{q}_k$ is linearly independent.
\end{proposition}

\begin{proof}
Suppose $c_1\vec q_1+\cdots+c_k\vec q_k=\vec 0$. Dot both sides with $\vec q_j$. By orthogonality
every term $\inner{\vec q_i}{\vec q_j}$ vanishes except $i=j$:
\[
  0 = \Big\langle \textstyle\sum_i c_i\vec q_i,\ \vec q_j\Big\rangle
    = c_j\inner{\vec q_j}{\vec q_j} = c_j\norm{\vec q_j}^2.
\]
Since $\vec q_j\neq\vec 0$, $\norm{\vec q_j}^2>0$, so $c_j=0$. This holds for each $j$, so all
coefficients are zero: the vectors are independent.
\end{proof}

If on top of being orthogonal each vector also has \emph{unit length}, the structure simplifies further.

\begin{definition}\label{ch17:def-orthonormal}
Vectors $\vec q_1,\dots,\vec q_k\in\R^n$ form an \textbf{orthonormal set} if they are mutually
orthogonal \emph{and} each is a unit vector:
\[
  \inner{\vec q_i}{\vec q_j} \;=\; \delta_{ij} \;=\;
  \begin{cases} 1 & i=j,\\ 0 & i\neq j.\end{cases}
\]
The symbol $\delta_{ij}$ is the \emph{Kronecker delta}: $1$ on the diagonal, $0$ off it.
\end{definition}

\begin{intuitionbox}
An orthonormal set consists of mutually perpendicular unit vectors, forming a rigid coordinate
frame in $\R^n$. The standard basis vectors $\vec{e}_1,\dots,\vec{e}_n$ form the simplest example,
but any rotation or reflection of this standard frame remains orthonormal. Because the vectors
are unit length and mutually orthogonal, coordinates decouple completely.
\end{intuitionbox}

Their practical advantage is immediate: in an orthonormal basis, finding coordinates does not
require solving a linear system; it requires only evaluating dot products. Suppose
$\vec q_1,\dots,\vec q_n$ is an orthonormal basis of $\R^n$ and you want to write a vector
$\vec v$ in this basis:
$\vec v = c_1\vec q_1+\cdots+c_n\vec q_n$.
Dotting both sides with $\vec q_j$, every term with $i \neq j$ vanishes:
\[
  \inner{\vec v}{\vec q_j} = \sum_i c_i\inner{\vec q_i}{\vec q_j} = c_j.
\]
So $c_j = \inner{\vec v}{\vec q_j}$ -- the $j$th coordinate is simply the dot product of $\vec v$
with the $j$th basis vector.

\begin{focusbox}
\textbf{Orthonormal expansion.} If $\vec q_1,\dots,\vec q_n$ is an orthonormal basis of $\R^n$,
then every $\vec v$ has the explicit expansion
\[
  \vec v \;=\; \inner{\vec v}{\vec q_1}\,\vec q_1 + \inner{\vec v}{\vec q_2}\,\vec q_2
            + \cdots + \inner{\vec v}{\vec q_n}\,\vec q_n .
\]
Each coefficient is a single dot product. This explicit formula avoids matrix inversion,
making orthonormal bases computationally stable.
\end{focusbox}

\begin{recoverybox}
\textbf{Pause and think: Direct coordinates.} Suppose $\vec{q}_1,\vec{q}_2,\vec{q}_3$ is an
orthonormal basis of $\R^3$. Without doing any elimination or writing down coordinates, what is
$\inner{5\vec{q}_1 - 4\vec{q}_2 + 7\vec{q}_3}{\vec{q}_2}$?
By orthonormality, $\inner{\vec{q}_1}{\vec{q}_2}=0$, $\inner{\vec{q}_3}{\vec{q}_2}=0$, and
$\inner{\vec{q}_2}{\vec{q}_2}=1$. The answer is $-4$. Each orthonormal basis vector projects
out only its corresponding component.
\end{recoverybox}

\begin{historybox}
Systematic orthogonalization grew out of curve-fitting and integral equations around the turn of the twentieth century. In 1883, the Danish mathematician and actuary J\o{}rgen Gram was calculating best-fit curves for forestry data. When fitting curves using ordinary powers ($1, t, t^2, \dots$), every extra degree scrambled all previous coefficients, requiring huge, unstable systems of equations to be resolved from scratch. Gram realized that replacing polynomials with mutually perpendicular functions caused the cross-terms to vanish, allowing each coefficient to be computed independently.

Two decades later, in 1907, Erhard Schmidt and his mentor David Hilbert at G\"ottingen were tackling \emph{linear integral equations}---equations where the unknown is an entire continuous curve trapped inside an integral, effectively a system with infinitely many variables. At the time, mathematicians possessed only classical calculus, determinants, and Fourier's trigonometric series; the concept of an abstract ``vector space'' did not yet exist. When Schmidt adapted Gram's orthogonalization procedure to continuous functions, the infinite-dimensional equations decoupled into independent one-dimensional projections, completely bypassing matrix inversion. Hilbert recognized the underlying geometry at work and formalized the framework of \emph{Hilbert space}---an infinite-dimensional arena where functions behave like vectors and integrals play the role of dot products.

At first, applying geometric words like ``perpendicular'' and ``projection'' to spaces of infinitely many dimensions startled classical analysts accustomed to explicit coordinate formulas. Within two decades, however, Werner Heisenberg, Max Born, and Erwin Schr\"odinger found that Hilbert's orthogonal function spaces provided the natural mathematical setting for quantum mechanics. Today, orthogonality underpins signal processing, computer graphics, and stable algorithms like the QR factorization.
\end{historybox}

\begin{example}\label{ch17:rotated-basis}
The rotated pair
\[
  \vec q_1=\tfrac{1}{\sqrt2}\colvec{1\\1},\qquad
  \vec q_2=\tfrac{1}{\sqrt2}\colvec{-1\\1}
\]
is orthonormal: $\inner{\vec q_1}{\vec q_1}=\tfrac12(1+1)=1$, likewise for $\vec q_2$, and
$\inner{\vec q_1}{\vec q_2}=\tfrac12(-1+1)=0$. To express $\vec v=\colvec{3\\1}$ in this basis,
just dot:
\[
  c_1=\inner{\vec v}{\vec q_1}=\tfrac{1}{\sqrt2}(3+1)=\tfrac{4}{\sqrt2}=2\sqrt2,\qquad
  c_2=\inner{\vec v}{\vec q_2}=\tfrac{1}{\sqrt2}(-3+1)=-\sqrt2.
\]
Sanity check: $2\sqrt2\,\vec q_1 - \sqrt2\,\vec q_2 = \colvec{2\\2}+\colvec{1\\-1}=\colvec{3\\1}
=\vec v.$ It works, and we never solved a system.
\end{example}

Figure~\ref{ch17:tilted-grid} shows what this looks like: the orthonormal pair is just a copy of the
standard graph paper rotated by $45^\circ$, and the coordinates $c_1,c_2$ are the readings of $\vec v$
along those tilted axes.

\begin{figure}[H]
\centering
\begin{tikzpicture}[scale=0.95,>=Stealth,line cap=round]
  % faint standard grid
  \draw[help lines,spacegray!30] (-0.6,-0.6) grid (3.4,3.4);
  \draw[->,spacegray] (-0.6,0)--(3.6,0) node[right]{$x$};
  \draw[->,spacegray] (0,-0.6)--(0,3.6) node[above]{$y$};
  % tilted orthonormal axes (q1 along (1,1)/sqrt2, q2 along (-1,1)/sqrt2)
  \draw[vecblue!55,thick,-{Stealth}] (0,0)--(2.55,2.55) node[vecblue!60!black,above right] {$\vec q_1$};
  \draw[vecred!55,thick,-{Stealth}] (0,0)--(-1.7,1.7) node[vecred!70!black,above left] {$\vec q_2$};
  % unit basis arrows
  \draw[->,vecblue,very thick] (0,0)--(0.707,0.707);
  \draw[->,vecred,very thick]  (0,0)--(-0.707,0.707);
  % vector v=(3,1)
  \draw[->,structurepurple,very thick] (0,0)--(3,1) node[above right=1pt]{$\vec v$};
  % drop to q1 axis: foot at c1*q1 = 2sqrt2*(.707,.707)=(2,2)
  \draw[spacegray!70,dashed] (3,1)--(2,2);
  \fill[black] (2,2) circle (1.0pt);
  \node[font=\small,below right] at (2.2,1.8) {$c_1\vec q_1$};
  \fill[black] (0,0) circle (1.1pt);
\end{tikzpicture}
\caption{The orthonormal basis $\vec q_1,\vec q_2$ is the standard grid turned $45^\circ$. The
coordinate $c_1=\inner{\vec v}{\vec q_1}$ is found by dropping $\vec v$ perpendicularly onto the
$\vec q_1$ axis (dashed line), landing at $c_1\vec q_1=(2,2)$; $c_2=\inner{\vec v}{\vec q_2}$ is the
matching reading along $\vec q_2$, landing at $c_2\vec q_2=(1,-1)$. Coordinates by dot product, no system to solve.}
\label{ch17:tilted-grid}
\end{figure}

\begin{connectionbox}
This expansion is the finite-dimensional ancestor of the \emph{Fourier series}. In a function
space, the sines and cosines $\{\sin(kx),\cos(kx)\}$ are mutually orthogonal under the inner
product $\inner{f}{g}=\int f(x)g(x)\,dx$, and the Fourier coefficients are computed by exactly the
same recipe: $c_k = \inner{f}{\,\text{that basis function}}$, an integral instead of a sum. Whole
fields -- signal processing, the heat equation, quantum mechanics -- are built on choosing a clever
orthonormal basis and reading off coordinates by integration. The chapter on Hermitian matrices and
Fourier makes this precise, but the idea is already in front of you: orthonormal means
``coordinates by dot product.''
\end{connectionbox}

\section{Orthogonal matrices and isometries}

Collect an orthonormal set into the columns of a matrix and $\mat Q\T\mat Q$ collapses to something
very simple.

\begin{definition}\label{ch17:def-orthogonal-matrix}
A square matrix $\mat Q$ is an \textbf{orthogonal matrix} if its columns form an orthonormal set,
equivalently if
\[
  \mat Q\T\mat Q \;=\; \mat I .
\]
\end{definition}

Let us see why those two statements are the same. If the columns are $\vec q_1,\dots,\vec q_n$,
then the $(i,j)$ entry of $\mat Q\T\mat Q$ is (row $i$ of $\mat Q\T$) times (column $j$ of $\mat Q$),
which is $\vec q_i\T\vec q_j=\inner{\vec q_i}{\vec q_j}$. Orthonormality says that equals
$\delta_{ij}$, and the matrix of $\delta_{ij}$ entries is precisely $\mat I$. So
``columns orthonormal'' and ``$\mat Q\T\mat Q=\mat I$'' say the same thing, entry by entry.

\begin{trapbox}
Despite the standard name, an ``orthogonal matrix'' has \emph{orthonormal} columns,
not merely orthogonal columns. Each column must have unit length. A matrix with
orthogonal columns of varying lengths is not an orthogonal matrix; the formal
condition is $\mat Q\T\mat Q = \mat I$.
\end{trapbox}

For a \emph{square} $\mat Q$, the equation $\mat Q\T\mat Q=\mat I$ says $\mat Q\T$ is a left inverse,
and since $\mat Q$ is square a left inverse is the inverse. So we get the cleanest inverse in
mathematics:
\[
  \boxed{\;\mat Q\inv = \mat Q\T\;}
\]
To invert an orthogonal matrix you do not solve anything -- you just transpose it. (As a bonus,
$\mat Q\T$ is then orthogonal too, so the \emph{rows} of an orthogonal matrix are also orthonormal,
and $\mat Q\mat Q\T=\mat I$ as well.)

\begin{example}\label{ch17:rotation-matrix}
The plane rotation by angle $\theta$,
\[
  \mat Q=\begin{bmatrix}\cos\theta & -\sin\theta\\ \sin\theta & \cos\theta\end{bmatrix},
\]
is orthogonal. Its columns each have length $\sqrt{\cos^2\theta+\sin^2\theta}=1$ and dot to
$-\cos\theta\sin\theta+\sin\theta\cos\theta=0$. Its inverse is the transpose, which is rotation by
$-\theta$ -- exactly what undoes a rotation. A reflection, e.g.\
$\begin{bmatrix}1&0\\0&-1\end{bmatrix}$ (flip across the $x$-axis), is orthogonal too. Rotations
and reflections are the orthogonal matrices of the plane.
\end{example}

\subsection{What orthogonal matrices preserve}

Here is the property that makes orthogonal matrices central to numerical linear algebra: they
move space \emph{rigidly}. They never stretch, squash, or shear. Lengths come out unchanged, and so
do angles.

\begin{theorem}[Orthogonal matrices preserve dot products]\label{ch17:Q-preserves}
If $\mat Q$ is orthogonal, then for all $\vec x,\vec y\in\R^n$,
\[
  \inner{\mat Q\vec x}{\mat Q\vec y} = \inner{\vec x}{\vec y}.
\]
In particular $\norm{\mat Q\vec x}=\norm{\vec x}$ (lengths preserved), and the angle between
$\mat Q\vec x$ and $\mat Q\vec y$ equals the angle between $\vec x$ and $\vec y$ (angles preserved).
\end{theorem}

\begin{proof}
Use the transpose rule $(\mat Q\vec x)\T=\vec x\T\mat Q\T$ and the defining identity
$\mat Q\T\mat Q=\mat I$:
\[
  \inner{\mat Q\vec x}{\mat Q\vec y}
  = (\mat Q\vec x)\T(\mat Q\vec y)
  = \vec x\T \mat Q\T\mat Q\,\vec y
  = \vec x\T\mat I\,\vec y
  = \vec x\T\vec y
  = \inner{\vec x}{\vec y}.
\]
Taking $\vec y=\vec x$ gives $\norm{\mat Q\vec x}^2=\norm{\vec x}^2$, hence equal lengths. Since
$\cos\theta=\inner{\vec x}{\vec y}/(\norm{\vec x}\norm{\vec y})$ and every quantity in that ratio is
preserved, the angle is preserved too.
\end{proof}

\begin{intuitionbox}
An orthogonal matrix is a \emph{rigid motion} of space that fixes the origin: a rotation, a
reflection, or a combination. Drop a constellation of points into $\R^n$, hit them all with
$\mat Q$, and the whole constellation turns or flips as one solid object -- every distance between
points and every angle is exactly what it was. Nothing is distorted. That is the geometric meaning
of $\mat Q\T\mat Q=\mat I$.
\end{intuitionbox}

\begin{figure}[H]
\centering
\begin{tikzpicture}[scale=1.05,>=Stealth]
  % left: unit square with a tilted vector
  \begin{scope}
    \draw[->,spacegray](-0.3,0)--(2.1,0); \draw[->,spacegray](0,-0.3)--(0,2.1);
    \draw[fill=sourceblue!12,draw=sourceblue] (0,0)--(1,0)--(1,1)--(0,1)--cycle;
    \draw[->,vecblue,very thick] (0,0)--(1,0);
    \draw[->,vecred,very thick] (0,0)--(0,1);
    \node at (0.5,-0.55) {before};
  \end{scope}
  % arrow
  \draw[->,very thick,structurepurple] (2.6,0.5)-- node[above]{$\mat Q$} (3.8,0.5);
  % right: rotated square -- same shape, turned
  \begin{scope}[shift={(5.6,0)}]
    \draw[->,spacegray](-0.9,0)--(1.6,0); \draw[->,spacegray](0,-0.3)--(0,2.1);
    \def\a{35}
    \draw[fill=sourceblue!12,draw=sourceblue,rotate=\a] (0,0)--(1,0)--(1,1)--(0,1)--cycle;
    \draw[->,vecblue,very thick,rotate=\a] (0,0)--(1,0);
    \draw[->,vecred,very thick,rotate=\a] (0,0)--(0,1);
    \node at (0.2,-0.55) {after};
  \end{scope}
\end{tikzpicture}
\caption{An orthogonal matrix turns (or flips) the unit square without distorting it. The sides
keep length $1$, the corner keeps its right angle: a rigid motion. Contrast a general matrix, which
would stretch the square into a slanted parallelogram.}
\label{ch17:rigid-motion}
\end{figure}

\begin{labbox}
Verify rigidity by hand with $\theta=90^\circ$, so $\mat Q=\begin{bmatrix}0&-1\\1&0\end{bmatrix}$.
Take $\vec x=\colvec{2\\1}$, $\vec y=\colvec{1\\3}$. Then $\mat Q\vec x=\colvec{-1\\2}$ and
$\mat Q\vec y=\colvec{-3\\1}$. Lengths: $\norm{\vec x}^2=5=\norm{\mat Q\vec x}^2$ and
$\norm{\vec y}^2=10=\norm{\mat Q\vec y}^2$. Dot products: $\inner{\vec x}{\vec y}=2+3=5$ and
$\inner{\mat Q\vec x}{\mat Q\vec y}=3+2=5$. Everything is preserved -- the $90^\circ$ turn moved the
arrows but kept all measurements.
\end{labbox}

\begin{remark}
Taking determinants of $\mat Q\T\mat Q=\mat I$ gives $\det(\mat Q)^2=1$, so $\det\mat Q=\pm1$.
The value $+1$ is a pure rotation (a ``proper'' motion that preserves handedness), and $-1$ is a
reflection or a rotation-times-reflection (it flips orientation, turning a left hand into a right
hand). Either way the volume scaling is $|\det\mat Q|=1$: orthogonal matrices preserve volume, as
a rigid motion must. We will meet determinants properly in a later chapter; for now, note how
cleanly that $\pm1$ falls out.
\end{remark}

\begin{recoverybox}
\textbf{Pause and think: Real eigenvalues of a rigid motion.} If an orthogonal matrix $\mat Q$ has a
real eigenvalue $\lambda$ with non-zero eigenvector $\vec x$, what could $\lambda$ possibly be?
Since $\mat Q$ preserves lengths, $\norm{\mat Q\vec x} = \norm{\lambda \vec x} = |\lambda|\,\norm{\vec x}$.
Equating this to $\norm{\vec x}$ forces $|\lambda|=1$. So real eigenvalues of an orthogonal matrix
can only be $+1$ (axes left untouched, like the rotation axis of a 3D spinning top) or $-1$
(directions flipped directly backwards, like in a mirror reflection). A rigid motion can never stretch
or shrink an eigenvector!
\end{recoverybox}

\begin{applicationbox}
\textbf{Why machine learning and graphics rely on $\mat Q\T\mat Q=\mat I$.}
\begin{itemize}
  \item \emph{Numerical stability.} Multiplying by $\mat Q$ never amplifies error, because it
        never amplifies length ($\norm{\mat Q\vec x}=\norm{\vec x}$). Algorithms like \emph{QR}
        factorization (decomposing a matrix into an orthogonal $\mat Q$ and upper-triangular $\mat R$)
        and the \emph{SVD} (singular value decomposition) are built from orthogonal matrices precisely
        so that rounding errors don't blow up. A general matrix can magnify a tiny error by a huge factor;
        an orthogonal one cannot.
  \item \emph{3D graphics and robotics.} A camera turn or a robot-arm joint is an orthogonal
        $3\times3$ matrix with $\det=+1$ (a rotation). Storing the inverse is free -- just transpose
        -- and repeated rotations never silently distort the model, because each one preserves
        every length and angle.
  \item \emph{Random projections and PCA.} Principal component analysis (PCA, which finds the dominant
        directions of a data cloud) uses eigenvectors that form an orthonormal basis; projecting data onto
        them rotates the cloud onto its natural axes without squashing distances. ``Orthonormal directions''
        is the data-science way of saying ``uncorrelated, unit-variance coordinates.''
  \item \emph{Whitening and decorrelation.} Because $\mat Q$ preserves inner products, it turns one
        orthonormal frame into another -- the exact operation behind rotating data into a basis
        where features no longer interfere.
\end{itemize}
The thread tying these together is Theorem~\ref{ch17:Q-preserves}: when your transformation
preserves length and angle, geometry survives the computation intact.
\end{applicationbox}

\begin{strangbox}
Gilbert Strang emphasizes that orthogonal matrices are numerically ideal
transformations to compute with, and that $\mat Q\inv=\mat Q\T$ is the central
property to remember. After Strang, you will see this same $\mat Q$ in three
central factorizations: $\mat A=\mat Q\mat R$ (Gram--Schmidt, next chapters),
$\mat A=\mat Q\Lambda\mat Q\T$ for symmetric matrices (the spectral theorem),
and $\mat A=\mat U\Sigma\mat V\T$ (the SVD), where $\mat U$ and $\mat V$ are
both orthogonal. In every case the orthogonal factors perform isometries (rigid
rotations and reflections) while a diagonal factor handles stretching.
\end{strangbox}

\section{Splitting a vector into perpendicular parts}

We close by applying the orthogonal-complement structure, which leads straight into the next
chapter on projections. The focus box of the complement section established that $\R^n = V\oplus V^\perp$
for \emph{any} subspace $V$: Theorem~\ref{ch17:dim-complement} shows the dimensions add to $n$,
while Proposition~\ref{ch17:trivial-intersection} forces $V\cap V^\perp=\{\vec0\}$.
The direct sum ``$\oplus$'' means every vector splits \emph{uniquely} into a piece in
$V$ and a piece in $V^\perp$.

\begin{theorem}[Orthogonal decomposition]\label{ch17:orthog-decomp}
Let $V\subseteq\R^n$ be a subspace. Every $\vec x\in\R^n$ can be written in exactly one way as
\[
  \vec x = \vec x_V + \vec x_{V^\perp}, \qquad \vec x_V\in V,\ \ \vec x_{V^\perp}\in V^\perp .
\]
\end{theorem}

\begin{proof}
\emph{Existence.} Choose an orthonormal basis $\vec q_1,\dots,\vec q_r$ of $V$ (one exists; building
it is the job of Gram--Schmidt, two chapters on). Set $\vec x_V=\sum_{i=1}^r\inner{\vec x}{\vec q_i}\vec q_i$
and $\vec x_{V^\perp}=\vec x-\vec x_V$. Clearly $\vec x_V\in V$. To see $\vec x_{V^\perp}\in V^\perp$,
dot it with any basis vector $\vec q_j$:
\[
  \inner{\vec x-\vec x_V}{\vec q_j} = \inner{\vec x}{\vec q_j} - \inner{\vec x}{\vec q_j} = 0,
\]
so $\vec x_{V^\perp}$ is orthogonal to every $\vec q_j$, hence to all of $V$.

\emph{Uniqueness.} If $\vec x=\vec a+\vec b=\vec a'+\vec b'$ are two such splittings with
$\vec a,\vec a'\in V$ and $\vec b,\vec b'\in V^\perp$, then $\vec a-\vec a'=\vec b'-\vec b$ lies in
both $V$ and $V^\perp$. By Proposition~\ref{ch17:trivial-intersection} the only such vector is
$\vec 0$, so $\vec a=\vec a'$ and $\vec b=\vec b'$.
\end{proof}

\begin{intuitionbox}
Given any subspace $V \subseteq \R^n$, every vector $\vec{x}$ decomposes uniquely
into a component $\vec{x}_V \in V$ and an orthogonal component $\vec{x}_{V^\perp} \in V^\perp$.
The vector $\vec{x}_V$ is the orthogonal \emph{projection} of $\vec{x}$ onto $V$,
while $\vec{x}_{V^\perp} = \vec{x} - \vec{x}_V$ is the orthogonal error vector.
In Chapter~18, this decomposition forms the basis for orthogonal projections and
least-squares regression.
\end{intuitionbox}

\begin{figure}[H]
\centering
\begin{tikzpicture}[scale=1.25,>=Stealth,line cap=round]
  % subspace V drawn as a slanted line through the origin
  \coordinate (O) at (0,0);
  \draw[targetgreen!70,very thick] (-1.6,-0.8) -- (2.4,1.2) node[below right]{$V$};
  % x and its two pieces
  \coordinate (X) at (1.1,1.8);
  \coordinate (XV) at (1.6,0.8);   % shadow on V (foot of perpendicular)
  \draw[->,vecblue,very thick] (O) -- (X) node[above]{$\vec x$};
  \draw[->,targetgreen!60!black,very thick] (O) -- (XV) node[below right]{$\vec x_V$};
  \draw[->,structurepurple,very thick] (XV) -- (X);
  \node[structurepurple,above right] at (1.35,1.3) {$\vec x_{V^{\perp}}$};
  % right-angle mark where the leftover meets V
  \draw[spacegray] ($(XV)+(-0.16,-0.08)$) -- ($(XV)+(-0.24,0.08)$) -- ($(XV)+(-0.08,0.16)$);
  \fill[black] (O) circle (1.1pt);
\end{tikzpicture}
\caption{The orthogonal decomposition $\vec x=\vec x_V+\vec x_{V^{\perp}}$. The projection
$\vec x_V$ lies in $V$; the residual $\vec x_{V^{\perp}}$ is perpendicular to $V$ (note the
right-angle mark). This decomposition is unique, providing the basis for projection and least squares.}
\label{ch17:decomp-pic}
\end{figure}

\begin{example}\label{ch17:split-example}
Split $\vec x=\colvec{2\\3}$ along the line $V=\spn\!\left(\colvec{1\\0}\right)$ (the $x$-axis). The
unit vector is $\vec q_1=\colvec{1\\0}$, so
\[
  \vec x_V = \inner{\vec x}{\vec q_1}\,\vec q_1 = 2\colvec{1\\0}=\colvec{2\\0},
  \qquad
  \vec x_{V^\perp} = \vec x-\vec x_V = \colvec{0\\3}.
\]
Indeed $\vec x_V$ is on the $x$-axis, $\vec x_{V^\perp}$ is on the $y$-axis ($=V^\perp$), they are
perpendicular, and they sum to $\vec x$. The projection on the axis is $\colvec{2\\0}$; the
orthogonal residual is $\colvec{0\\3}$.
\end{example}

\begin{focusbox}
\textbf{Summary of core concepts.} Orthogonal vectors have zero dot product.
A subspace $V$ has a unique orthogonal complement $V^\perp$ such that $\dim V+\dim V^\perp=n$.
The four fundamental subspaces form two orthogonal complement pairs:
$\nullsp(\mat A)=\row(\mat A)^\perp$ in $\R^n$ and $\nullsp(\mat A\T)=\col(\mat A)^\perp$ in $\R^m$.
Orthonormal vectors allow coordinates to be calculated directly by dot products;
arranged into a square matrix, they form an orthogonal matrix $\mat Q$ with $\mat Q\inv=\mat Q\T$
that preserves lengths and angles. Every vector decomposes uniquely into an orthogonal sum
$\vec{x} = \vec{x}_V + \vec{x}_{V^\perp}$.
\end{focusbox}

\begin{connectionbox}
Orthonormality provides direct coordinates without matrix inversion. Finding coordinates relative to an orthonormal frame requires only dot products, and orthogonal matrices $\mat{Q}$ act as rigid motions (isometries, transformations that preserve lengths, angles, and distances exactly) without amplifying numerical errors.

Yet real-world data and columns of matrices are almost never orthogonal when they arrive. When a system $\mat{A}\vec{x}=\vec{b}$ is inconsistent because $\vec{b} \notin \col(\mat{A})$, knowing that $\col(\mat{A})$ and $\nullsp(\mat{A}\T)$ are orthogonal complements is not enough; we need an explicit operator to drop $\vec{b}$ onto $\col(\mat{A})$. In Chapter~18, we develop the algebraic machinery of \emph{projections}, building the projection matrix $\mat{P} = \mat{A}(\mat{A}\T\mat{A})\inv\mat{A}\T$ that finds the closest point on any subspace.
\end{connectionbox}

\section{Exercises}

\begin{exercise}\label{ch17:ex-basic-dot}
Decide which of these pairs are orthogonal, by computing the dot product:
(a) $\colvec{1\\-2\\1}$ and $\colvec{2\\1\\0}$;\quad
(b) $\colvec{3\\0\\-1}$ and $\colvec{1\\5\\3}$;\quad
(c) $\colvec{1\\1\\1\\1}$ and $\colvec{1\\1\\-1\\-1}$.
For each orthogonal pair, confirm Pythagoras: $\norm{\vec u-\vec v}^2=\norm{\vec u}^2+\norm{\vec v}^2$.
\end{exercise}

\begin{exercise}\label{ch17:ex-complement-line}
Let $V=\spn\!\left(\colvec{2\\-1\\2}\right)\subseteq\R^3$. Find a basis for $V^\perp$ by solving the
single equation $2x-y+2z=0$. Confirm that $\dim V+\dim V^\perp=3$, and verify each of your basis
vectors really is orthogonal to $\colvec{2\\-1\\2}$.
\end{exercise}

\begin{exercise}\label{ch17:ex-complement-plane}
Let $W=\spn\!\left(\colvec{1\\0\\1},\colvec{0\\1\\1}\right)\subseteq\R^3$. Find $W^\perp$.
(Hint: stack the two vectors as rows of a $2\times3$ matrix $\mat A$; then
$W^\perp=\row(\mat A)^\perp=\nullsp(\mat A)$.) What is $\dim W^\perp$, and does it match
Theorem~\ref{ch17:dim-complement}?
\end{exercise}

\begin{exercise}\label{ch17:ex-four-spaces}
For $\mat A=\begin{bmatrix}1&2&3\\2&4&6\end{bmatrix}$, find the rank and a basis for each of the four
fundamental subspaces. Then explicitly check the two orthogonality relations of
Theorem~\ref{ch17:two-pairs}: nullspace $\perp$ row space, and left nullspace $\perp$ column space.
\end{exercise}

\begin{exercise}\label{ch17:ex-orthonormal-coords}
The vectors $\vec q_1=\tfrac15\colvec{3\\4}$ and $\vec q_2=\tfrac15\colvec{-4\\3}$ are orthonormal
(check it). Use the orthonormal expansion formula to write $\vec v=\colvec{5\\10}$ as
$c_1\vec q_1+c_2\vec q_2$, and verify your answer by recombining.
\end{exercise}

\begin{exercise}\label{ch17:ex-build-Q}
Build a $3\times3$ orthogonal matrix whose first column is $\tfrac13\colvec{2\\2\\1}$.
(Find a second unit vector orthogonal to it, then a third orthogonal to both -- a cross product is a
fast way to get the third.) Verify $\mat Q\T\mat Q=\mat I$ for your matrix.
\end{exercise}

\begin{exercise}\label{ch17:ex-preserve-proof}
Prove directly from $\mat Q\T\mat Q=\mat I$ that the columns of an orthogonal matrix are
orthonormal, and conversely. Then prove that if $\mat Q_1$ and $\mat Q_2$ are both orthogonal
$n\times n$ matrices, so is their product $\mat Q_1\mat Q_2$. (Geometrically: composing two rigid
motions gives a rigid motion.)
\end{exercise}

\begin{exercise}\label{ch17:ex-split}
Let $V=\spn\!\left(\colvec{1\\1\\0}\right)$ and $\vec x=\colvec{3\\1\\4}$. Compute the orthogonal
decomposition $\vec x=\vec x_V+\vec x_{V^\perp}$ from Theorem~\ref{ch17:orthog-decomp}.
(First normalize the spanning vector to a unit $\vec q_1$.) Confirm that
$\inner{\vec x_V}{\vec x_{V^\perp}}=0$ and that the two pieces add back to $\vec x$.
\end{exercise}

\begin{exercise}\label{ch17:ex-perp-perp-proof}
Prove that for any two subspaces, $V\subseteq W$ implies $W^\perp\subseteq V^\perp$ (taking
complements reverses inclusion). Use this to give a short argument that $((V^\perp)^\perp)^\perp=V^\perp$.
\end{exercise}

\begin{exercise}\label{ch17:ex-challenge-rotation}
\textbf{(Challenge.)} Show that \emph{every} $2\times2$ orthogonal matrix is either
$\begin{bmatrix}\cos\theta&-\sin\theta\\\sin\theta&\cos\theta\end{bmatrix}$ (a rotation, $\det=+1$)
or $\begin{bmatrix}\cos\theta&\sin\theta\\\sin\theta&-\cos\theta\end{bmatrix}$ (a reflection,
$\det=-1$) for some angle $\theta$. (Hint: the first column is a unit vector, so it equals
$\colvec{\cos\theta\\\sin\theta}$ for some $\theta$; the second column is a unit vector orthogonal
to it, leaving exactly two choices.) Interpret each case geometrically.
\end{exercise}

\begin{exercise}\label{ch17:ex-proof-distance}
\textbf{(Proof.)} Use the orthogonal decomposition $\vec x=\vec x_V+\vec x_{V^\perp}$ to prove the
Pythagorean identity $\norm{\vec x}^2=\norm{\vec x_V}^2+\norm{\vec x_{V^\perp}}^2$. Then explain in
one sentence why this means the projection $\vec x_V$ is \emph{never longer} than $\vec x$ itself --
a fact you will use constantly in the projections chapter.
\end{exercise}

\begin{exercise}\label{ch17:ex-fredholm-solvability}
\textbf{(Solvability and the left nullspace.)}
Theorem~\ref{ch17:two-pairs} established that $\nullsp(\mat A\T) = \col(\mat A)^\perp$, which by
Corollary~\ref{ch17:perp-perp} means $\col(\mat A) = (\nullsp(\mat A\T))^\perp$. Explain in two
plain sentences why this provides the universal test for solvability: the linear system
$\mat A\vec x = \vec b$ has an exact solution if and only if $\vec b$ is orthogonal to every vector
in the left nullspace $\nullsp(\mat A\T)$. (This celebrated result is the finite-dimensional
\emph{Fredholm alternative}.)
\end{exercise}
